\section{CREATE DATABASE dan USE}

Sebelum menyimpan data, langkah pertama adalah membuat \textbf{database} --- wadah logis yang menampung kumpulan tabel terkait. Perintah \code{CREATE DATABASE nama\_database;} membuat database baru; tambahkan \code{IF NOT EXISTS} agar tidak error jika database sudah ada. Konvensi penamaan: gunakan huruf kecil, tanpa spasi (gunakan underscore), dan nama yang deskriptif --- misalnya \code{toko\_online}, \code{sistem\_akademik}, atau \code{blog\_pribadi}. Setelah dibuat, database dapat dilihat dengan \code{SHOW DATABASES;} \cite{mysqlDocs}.

Untuk bekerja dengan database tertentu, gunakan perintah \code{USE nama\_database;}. Semua perintah SQL berikutnya (CREATE TABLE, SELECT, INSERT, dll.) akan beroperasi pada database yang sedang aktif. Untuk menghapus database (beserta seluruh tabel dan datanya), gunakan \code{DROP DATABASE nama\_database;} --- perintah ini \textbf{permanen dan tidak dapat di-undo}, sehingga harus digunakan dengan sangat hati-hati \cite{mysqlGettingStarted}.

\begin{webcode}[language=SQL, caption={Membuat dan menggunakan database}]
-- Membuat database dengan charset UTF-8
CREATE DATABASE IF NOT EXISTS toko_online
  CHARACTER SET utf8mb4
  COLLATE utf8mb4_unicode_ci;

-- Menggunakan database
USE toko_online;

-- Melihat database yang sedang aktif
SELECT DATABASE();

-- Menghapus database (HATI-HATI!)
-- DROP DATABASE toko_online;
\end{webcode}

\begin{catatan}
Selalu gunakan charset \code{utf8mb4} dan collation \code{utf8mb4\_unicode\_ci} saat membuat database baru. Charset \code{utf8mb4} mendukung seluruh karakter Unicode termasuk emoji, sedangkan \code{utf8} bawaan MySQL hanya mendukung 3 byte (tidak lengkap). Pengaturan ini dapat dilakukan saat CREATE DATABASE atau di level server melalui \code{my.cnf}/\code{my.ini}.
\end{catatan}

